\section{Related Work}

\subsection{\LaTeX\ usability}

\LaTeX\ affords substantial control over technical and scientific documents.
While experienced \LaTeX\ users may be more efficient in drafting their documents, the same is not observed among first-time or one-time users.
\LaTeX\ is good at rendering mathematical expressions \cite{knauff_efficiency_2014}.
General-purpose document processing softwares like Microsoft Word, Google Docs, and LibreOffice recognise this and allow authors to embed snippets of \LaTeX\ to render mathematical sections within their documents \cite{matthews_craft_2019}.
On the other hand, rendering a table is difficult to get right due to its complicated syntax, and even experts regularly struggle with this task \cite{knauff_efficiency_2014}.

\subsection{Automated \LaTeX\ generation}

The complexity of \LaTeX\ syntax can be attributed to the extent of detail and control it affords.
Unfortunately, the consequence is that rendering even a basic component can become difficult.
For this reason, researchers have explored ways to simplify the process of writing \LaTeX\ code.
To address the challenge of creating tables in \LaTeX, Kayal et al. built a system to recognise tables using OCR and render tables using a ResNet transformer \cite{kayal_tables_2023}.
A similar solution was explored by Xia et al. to convert images and tables to \LaTeX \cite{xia_latexnet_2025}.
They collected images and tables along with their corresponding \LaTeX\ source code from the preprint repository ArXiv and trained an image-multiplexer convolutional neural network.
They observed high similarity between the generated and expected code.
However, these solutions are developed as prototypes.
Also, as they are trained on data instead of the syntax, they fail at edge cases and unexpected requests.

OpenAI developed a tool called Prism that closely resembles the functionality of Spartan Write \cite{openai_prism}.
Prism allows researchers to edit their documents using natural language.
The tool maintains the \LaTeX\ code for the document internally, similar to Spartan Write's architecture.
However, Prism is more relaxed in its vision of academic integrity.
It performs proofreading, edits, literature surveys, and other operations that freely alter the document contents.
As discussed in the following section, academia is generally opposed to the use of generated content in research publications.
Spartan Write aims to find a balance, bringing the benefits of AI assistance while upholding this.

Tools similar to Prism are in active development.
An open-source alternative called OpenPrism \cite{github_openprism} seeks to rebuild Prism's functionality using open-source tools.
Bibby AI \cite{bibby_ai} is a tool that generates code from natural language requests.
Autype is another tool that eliminates the usage of \LaTeX\ code, instead opting for markdown-like syntax \cite{autype}.
Being imitations of Prism, they share the same issues that Prism faces.

\subsection{Generative AI in academic writing}

The surge in LLMs and GenAI has seen large-scale adoption in disciplines like creative work and software development.
However, students and researchers in academia are strongly advised against using GenAI in their work.
Studies show increased usage of LLMs across publications and associate it with lowered material quality, communication inaccuracy, abnormal patterns, and other issues \cite{kobak_delving_2025}.
For these reasons, surveys reveal that GenAI usage is generally shunned in academia \cite{kwon_survey_2025}.

A review summarizes eleven studies between 2021 and 2023 on the usage of AI in academic writing.
It reports productivity wins in tasks like plagiarism detection and language learning support \cite{guo_review_2024}.
However, the studies also emphasised concerns regarding integrity, accuracy, plagiarism, and over-reliance.

Another survey compiled information from 24 studies and identified various ways in which AI positively impacts the academic writing process.
The main factors were generating ideas, improvng content structure, conducting a literature review, managing and analysing data, and proofreading.
They, too, cite concerns regarding integrity.

PaperDebugger \cite{paperdebugger} is an extension system developed as an overhaul to Overleaf, the popular online LaTeX editor.
It uses a plugin-based multi-agent system to assist with common editing tasks.
This tool is widely used as evidenced by the download metrics in the repository.

The general consensus is that AI is extremely helpful in the process.
If it is deployed in a way that does not breach the ethical principles of the academic community, it offers considerable benefits to authors and publishers.

\subsection{Other tools}

Surveying other tools that may be able to solve the same problem as Spartan Write, it is possible to use agentic harnesses like Cursor, Copilot, Claude, or OpenAI Codex to perform the same tasks.
However, Spartan Write differs from these solutions in several ways.
First, it is a purpose-built tool for academic writing rather than a general-purpose tool.
Safeguards for integrity can be handcrafted based on iteration.
The templates are curated by hand to guarantee consistent outputs.
The agent is coupled with the compiler and everything ships as one package.
This allows the agent to be aware of failures or warnings generated during compilation.
Lastly, using a dedicated framework enables shipping updates with fixes and new features.
